\section{The formalization in Lean}\label{app:lean}

The numbered statements of Sections~\ref{sec:setting}--\ref{sec:hollowcore} and Appendix~\ref{app:LR} --- the definitions,
lemmas, propositions, theorems and corollaries listed in the table below --- have been formalized and proved in the Lean~4
proof assistant against the Mathlib library, with the exceptions of the exact count of Theorem~\ref{thm:kernel} --- its lower
bound, the existence of the zero modes constructed in Appendix~\ref{app:kernel}, is proved, and it is all the main theorems use
--- and Lemma~\ref{lem:units}, and, in Section~\ref{sec:index}, the restriction to $U=\Theta$ described in (a).\footnote{Lean 4.34.1 and Mathlib
at commit d13f23b7. The formalization is the Lean project \texttt{HRTorus} in the directory \texttt{lean/} of
\url{https://github.com/bbradlyn/hr-torus-proof}; its \texttt{README.md} gives the build, the axiom checks and the kernel
replay, and \texttt{lean/MAP.md} the complete correspondence: declaration by declaration for the modules written in the
formalization's stages 1 and F, by statement for the first commit's.} The project has forty-three
modules and about 50{,}100 lines. No proof is left open: every declaration depends only on the three standard axioms of
Lean's logic (propositional extensionality, choice and quotient soundness), and every module has been replayed through Lean's
kernel by \texttt{leanchecker}, independently of the elaborator. What the kernel certifies is the Lean text. Whether that text says what this paper says is
a matter of reading, and the table is the dictionary for it.

\paragraph{What is proved.} Three kinds of statement should be distinguished.
\begin{enumerate}[(a)]
\item \emph{Unconditional.} Everything in Sections~\ref{sec:setting}--\ref{sec:hollowcore} and Appendix~\ref{app:LR}
except the exact count of Theorem~\ref{thm:kernel}, Lemma~\ref{lem:units} (below) and the general-$U$ clauses of
Section~\ref{sec:index} named next; and, of Appendix~\ref{app:bbdf}, Lemmas~\ref{lem:bbdf41}, \ref{lem:expbounds}
and~\ref{lem:det}. In particular Theorem~\ref{thm:main}, Corollary~\ref{cor:oddN} and Corollary~\ref{cor:uniform} are
proved with no hypothesis on the zero modes (\texttt{hollow\_core\_index\_unconditional},
\texttt{hollow\_core\_E11\_tendsto\_unconditional}, \texttt{odd\_E3\_tendsto\_unconditional},
\texttt{odd\_many\_small\_unconditional}, \texttt{uniform\_E11\_tendsto\_unconditional},
\texttt{uniform\_E11\_eventually\_unconditional}); the divisibility half of Theorem~\ref{thm:main}(i), $4\mid p$, needs no
information about the zero modes at all (\texttt{hollow\_core\_index\_four\_dvd}), and the rest uses only the lower bound of
Theorem~\ref{thm:kernel} --- at least ten zero modes of $H_N$ in the sector $(N/2,N/2)$ at every even $N$, on every rectangular
torus, and at least two in each of the sectors $((N\pm1)/2,(N\mp1)/2)$ at every odd $N\ge3$ --- which is proved from the explicit
states of Proposition~\ref{kz:prop:explicit} through Lemma~\ref{kz:lem:A} in the direction (2)~$\Rightarrow$~(3)
(\texttt{ten\_le\_finrank\_ker}, \texttt{two\_le\_finrank\_ker\_odd}). Theorem~\ref{thm:index} is formalized for $U=\Theta$, the case this paper's numbered results use, for every
Hamiltonian of Definition~\ref{def:class} with (H1)--(H2) (\texttt{index\_theorem\_\ensuremath{\Theta}}), the inputs supplied with an
explicit $\varepsilon(N)=O(N^{-\infty})$ and $M(N)=(1+4\pi C_h\norm{W_\gamma}_1)N^2$; for a general $U$ only its finite-dimensional
core is formalized, in the following quantitative form: for every finite-dimensional complex inner product space and every
choice of the data of Setting~\ref{set:index}, Assumptions~\ref{ass:gap}--\ref{ass:U} (with $\gamma>\Delta$) and
Definitions~\ref{def:dressing} and~\ref{def:bbdfobjects} satisfying the exact relations stated there, those of
Lemma~\ref{lem:split}(i) and of the filters (Definition~\ref{def:filters}, Lemma~\ref{lem:Wexists}) included, if each of the
twenty-seven quantities --- the twenty-one separated commutators of Lemma~\ref{lem:commutation}(ii), (U2), the $+$ half of
(U4) and the four clustering commutators of Corollary~\ref{cor:Zcommute} --- is at most $\varepsilon$ and
$\norm{T_-},\norm{K_-}\le M$, then there is an integer $n$ with $\abs{\Tr(PT_-)-n}\le400\,p\,(M+1)\sqrt{\varepsilon}$
(\texttt{index\_core\_400}). The general-$U$ clauses of Lemma~\ref{lem:objects} and Lemma~\ref{lem:commutation}(ii) are not formalized.
\item \emph{In two forms.} The second half of Theorem~\ref{thm:main}(i), $12\le p$, Theorem~\ref{thm:main}(ii),
Corollary~\ref{cor:oddN} and Corollary~\ref{cor:uniform} are also in the tree in the form in which their statements were first
reviewed, as implications from the named hypotheses \texttt{HasTenZeroModes} and \texttt{HasTwoZeroModes} --- exactly ten zero
modes in the sector $(N/2,N/2)$ at every even $N$ (for Corollary~\ref{cor:uniform}, on every torus of the family), respectively
exactly two in each of the sectors $((N\pm1)/2,(N\mp1)/2)$ at every odd $N\ge3$, the content of Theorem~\ref{thm:kernel}
together with $H_N\ge0$ (\texttt{hollow\_core\_index}, \texttt{hollow\_core\_E11\_tendsto}, \texttt{odd\_E3\_tendsto},
\texttt{odd\_many\_small}, \texttt{uniform\_E11\_tendsto}, \texttt{uniform\_E11\_eventually}). The unconditional forms are
derived from the same statements under the weaker hypotheses \texttt{HasAtLeastTenZeroModes} and
\texttt{HasAtLeastTwoZeroModes} (at least ten, at least two), which the lower bound discharges.
\item \emph{Not formalized.} The exact count of Theorem~\ref{thm:kernel}: in Appendix~\ref{app:kernel}, Lemmas~\ref{kz:lem:A1},
\ref{kz:lem:A2} and~\ref{kz:lem:A4}, the directions (1)~$\Leftrightarrow$~(2) of Lemma~\ref{kz:lem:A}, Lemma~\ref{kz:lem:theta}(b),
Lemma~\ref{kz:lem:B}(a), (c) and its surjectivity, Lemmas~\ref{kz:lem:C} and~\ref{kz:lem:D}, Propositions~\ref{kz:prop:induct}
and~\ref{kz:prop:base} and Corollary~\ref{kz:cor:bounds}, whose arguments rest on the argument principle and on the analysis of
several complex variables, neither of which the Mathlib version used provides; Appendix~\ref{app:twist} (Theorem~\ref{thm:twist}
and Corollary~\ref{cor:alltori}), which no Lean statement uses; and Lemma~\ref{lem:units}, the planar statement about the units of $V_1$, on which nothing depends: no result of
this paper uses it and no Lean declaration does; it fixes the units in which the gaps are measured.
\end{enumerate}

\paragraph{The correspondence.} The table lists every numbered statement of Sections~\ref{sec:setting}--\ref{sec:hollowcore} and
Appendices~\ref{app:LR} and~\ref{app:kernel} with the Lean module and the declarations that formalize it; a statement with several parts is
usually several declarations. The lemmas of Section~\ref{sec:index} that the finite-dimensional index core proves
(Lemmas~\ref{lem:bbdf42}--\ref{lem:determinant} and Theorem~\ref{thm:index} itself) are formalized in the twelve modules
\texttt{IndexCore.Statement} to \texttt{IndexCore.Sanity}, with \texttt{IndexCore.Translation} and \texttt{IndexCore.Ring} for
Section~\ref{sec:translation}; their declarations cite the statements, and the table names the modules and the main
declarations. Lemma~\ref{lem:frame-consistency}(i)--(ii)
and Lemma~\ref{lem:localization} are formalized in \texttt{SlotFrame.SlotFrame}, also of the first commit, and the table names its
declarations. Where more than four declarations formalize one statement the table
names four and counts the rest; \texttt{lean/MAP.md} names them all. The Lean statements follow this paper's text under the
conventions listed after the table.

% generated by lean/gen-map.py from papers/hr/main.aux and the Lean sources; do not edit by hand
\begin{footnotesize}
\begin{longtable}{@{}p{2.35cm}p{3.35cm}>{\raggedright\arraybackslash}p{9.6cm}@{}}
\toprule Statement & Title & Lean module: declarations \\ \midrule \endfirsthead
\toprule Statement & Title & Lean module: declarations \\ \midrule \endhead
\bottomrule \endlastfoot
Definition~2.1 & the torus and its lowest Landau level & \texttt{LLL.Space}: \texttt{Space}, \texttt{IsLLL}, \texttt{l2inner} \\
Lemma~2.2 & Landau-gauge orbitals & \texttt{LLL.Space}: \texttt{orbBasis}, \texttt{finrank\_\allowbreak{}space}, \texttt{psi\_\allowbreak{}isLLL}, \texttt{l2inner\_\allowbreak{}psi} and 5 more \\
Lemma~2.3 & evaluation vectors & \texttt{LLL.Space}: \texttt{ev}, \texttt{evD}, \texttt{inner\_\allowbreak{}ev}, \texttt{inner\_\allowbreak{}evD} and 9 more \\
Lemma~2.4 & magnetic translations & \texttt{LLL.Space}: \texttt{thxFun\_\allowbreak{}psi}, \texttt{thyFun\_\allowbreak{}psi}, \texttt{toFun\_\allowbreak{}thx}, \texttt{toFun\_\allowbreak{}thy} and 11 more \\
Definition~2.5 & Fock space, sectors, the many-body translation & \texttt{Fermion.Pilot.JW}: \texttt{IsEven}; \texttt{Fock.Physical}: \texttt{FockSpace}, \texttt{cf}, \texttt{cjs}, \texttt{njs} and 16 more; \texttt{LLL.Space}: \texttt{Space2} \\
Lemma~2.6 & sectors and the many-body translation & \texttt{Fock.Physical}: \texttt{orthogonalFamily\_\allowbreak{}sector}, \texttt{isInternal\_\allowbreak{}sector}, \texttt{Psec\_\allowbreak{}eq\_\allowbreak{}cfc}, \texttt{isSectorPreserving\_\allowbreak{}iff\_\allowbreak{}sum} and 11 more \\
Definition~2.7 & the hollow-core Hamiltonian & \texttt{Fock.Hamiltonian}: \texttt{Vcoef}, \texttt{HN}, \texttt{Vcoef\_\allowbreak{}latticeSum} \\
Lemma~2.8 & $V_1$ units & \emph{not formalized: the planar units statement; nothing in the formalization depends on it} \\
Theorem~2.10 & the zero modes & \texttt{Assembly.Main}: \texttt{hasTenZeroModes\_\allowbreak{}of\_\allowbreak{}kernel}, \texttt{hasTwoZeroModes\_\allowbreak{}of\_\allowbreak{}kernel} \emph{(assumed; the two bridges from the kernel dimensions to the named hypotheses)} \\
Definition~3.1 & admissible tori and the slot frame & \texttt{Fock.Slot}: \texttt{slotDist}; \texttt{LLL.Frame}: \texttt{frameVec}, \texttt{analysis}, \texttt{analysis2}; \texttt{Locality.Basic}: \texttt{xs}; \texttt{Locality.LiebRobinson}: \texttt{slotSetDist} \\
Lemma~3.2 & consistency of the construction & \texttt{Locality.Basic}: \texttt{card\_\allowbreak{}Ia\_\allowbreak{}lt}, \texttt{eq\_\allowbreak{}of\_\allowbreak{}mem\_\allowbreak{}Ia\_\allowbreak{}of\_\allowbreak{}intCast\_\allowbreak{}eq}, \texttt{wa\_\allowbreak{}add\_\allowbreak{}nx\_\allowbreak{}add\_\allowbreak{}N\ensuremath{\varphi}}; \texttt{SlotFrame.SlotFrame}: \texttt{window\_\allowbreak{}contDiff}, \texttt{window\_\allowbreak{}eq\_\allowbreak{}zero}, \texttt{window\_\allowbreak{}sq\_\allowbreak{}tsum} (i); \texttt{SlotFrame.SlotFrame}: \texttt{consistency\_\allowbreak{}ii}, \texttt{Ia\_\allowbreak{}nonempty} (ii) \\
Proposition~3.3 & Parseval & \texttt{LLL.Frame}: \texttt{sum\_\allowbreak{}rankOne\_\allowbreak{}frameVec}, \texttt{adjoint\_\allowbreak{}analysis\_\allowbreak{}comp\_\allowbreak{}analysis}, \texttt{norm\_\allowbreak{}analysis}, \texttt{analysis\_\allowbreak{}comp\_\allowbreak{}adjoint} and 5 more \\
Proposition~3.4 & slot geometry & \texttt{Locality.Basic}: \texttt{slotBall}, \texttt{strip}, \texttt{Cvol}, \texttt{Cstr} and 6 more \\
Lemma~3.5 & uniform localization & \texttt{LLL.Frame}: \texttt{slotProj}, \texttt{localization\_\allowbreak{}iii\_\allowbreak{}analysis}; \texttt{SlotFrame.SlotFrame}: \texttt{localization\_\allowbreak{}i}, \texttt{localization\_\allowbreak{}ii}, \texttt{localization\_\allowbreak{}iii}, \texttt{localization} \\
Lemma~3.6 & exact column support & \texttt{Locality.Basic}: \texttt{col}, \texttt{inner\_\allowbreak{}orb\_\allowbreak{}frameVec}, \texttt{inner\_\allowbreak{}orb\_\allowbreak{}frameVec\_\allowbreak{}eq\_\allowbreak{}zero}, \texttt{analysis\_\allowbreak{}orb\_\allowbreak{}mem\_\allowbreak{}col} and 2 more \\
Definition~3.7 & slot Fock space, local algebras, the embedding & \texttt{Fermion.Pilot.JW}: \texttt{chat}, \texttt{numberOp}, \texttt{CAR}, \texttt{Alg}; \texttt{Fock.Embedding}: \texttt{iota}, \texttt{iota\_\allowbreak{}cf}, \texttt{exists\_\allowbreak{}starAlgHom\_\allowbreak{}iota}; \texttt{Fock.Slot}: \texttt{chatVec}, \texttt{chatVec\_\allowbreak{}spinTensor\_\allowbreak{}single}, \texttt{SlotFock}, \texttt{fat} \\
Lemma~3.8 & properties of the embedding & \texttt{Fock.Embedding}: \texttt{FfA}, \texttt{FfA\_\allowbreak{}conjTranspose\_\allowbreak{}mul\_\allowbreak{}self}, \texttt{norm\_\allowbreak{}FfA\_\allowbreak{}apply}, \texttt{range\_\allowbreak{}FfA} and 14 more; \texttt{Fock.Slot}: \texttt{gauge\_\allowbreak{}image\_\allowbreak{}CAR}, \texttt{gauge\_\allowbreak{}image\_\allowbreak{}Alg} \\
Definition~3.9 & strips, columns & \texttt{Locality.Basic}: \texttt{colBand}, \texttt{cthickening\_\allowbreak{}circImage\_\allowbreak{}Ioo} \\
Lemma~3.10 & strips and fattenings & \texttt{Locality.Basic}: \texttt{fat\_\allowbreak{}strip\_\allowbreak{}subset}, \texttt{strip\_\allowbreak{}cthickening\_\allowbreak{}subset} \\
Definition~3.11 & almost supported operators & \texttt{Locality.Basic}: \texttt{TorusSeq}, \texttt{setNorm}, \texttt{ASBound}, \texttt{ASWitness} and 4 more \\
Definition~4.1 & triplet pairs, local terms, centre pair operators & \texttt{Hamiltonian.Pairs}: \texttt{chi}, \texttt{cpair}, \texttt{cpairHat}, \texttt{B\ensuremath{\nu}} and 7 more \\
Theorem~4.2 & the exact local form & \texttt{Hamiltonian.LocalForm}: \texttt{localForm\_\allowbreak{}a}, \texttt{localForm\_\allowbreak{}b}, \texttt{localForm\_\allowbreak{}c} \\
Lemma~4.3 & jets on the diagonal & \texttt{Hamiltonian.Pairs}: \texttt{D1x}, \texttt{D1}, \texttt{jetF}, \texttt{jets\_\allowbreak{}symm} and 1 more \\
Proposition~4.4 & uniform size and rapid locality & \texttt{Hamiltonian.SlotLocal}: \texttt{ChHC}, \texttt{B\ensuremath{\nu}R}, \texttt{hRr}, \texttt{norm\_\allowbreak{}B\ensuremath{\nu}\_\allowbreak{}le} and 6 more \\
Lemma~4.5 & covariance & \texttt{Hamiltonian.LocalForm}: \texttt{commute\_\allowbreak{}\ensuremath{\Theta}\_\allowbreak{}HN}; \texttt{Hamiltonian.Pairs}: \texttt{HD}, \texttt{\ensuremath{\Theta}\_\allowbreak{}conj\_\allowbreak{}B\ensuremath{\nu}}, \texttt{\ensuremath{\Theta}\_\allowbreak{}conj\_\allowbreak{}hR}, \texttt{\ensuremath{\Theta}\_\allowbreak{}conj\_\allowbreak{}HD} \\
Definition~4.6 & slot-local covariant Hamiltonians & \texttt{Hamiltonian.SlotLocal}: \texttt{H1Data}, \texttt{H1Data.H2}, \texttt{SatisfiesH1}, \texttt{SatisfiesH1H2} \\
Corollary~4.8 & the hollow-core Hamiltonian is in the class & \texttt{Hamiltonian.SlotLocal}: \texttt{hollowCore\_\allowbreak{}mem\_\allowbreak{}class} \\
Definition~4.9 & band charges & \texttt{Hamiltonian.BandCharges}: \texttt{Qband} \\
Lemma~4.10 & elementary properties of band charges & \texttt{Hamiltonian.BandCharges}: \texttt{commute\_\allowbreak{}Qband}, \texttt{Qband\_\allowbreak{}isSectorPreserving}, \texttt{spectrum\_\allowbreak{}Qband\_\allowbreak{}subset}, \texttt{norm\_\allowbreak{}Qband\_\allowbreak{}le} and 4 more \\
Lemma~4.11 & band charges in the slot lattice & \texttt{Hamiltonian.BandCharges}: \texttt{vj}, \texttt{bandX}, \texttt{iota\_\allowbreak{}njs\_\allowbreak{}up}, \texttt{vj\_\allowbreak{}support} and 9 more \\
Lemma~5.1 & conditional expectations & \texttt{Fermion.Pilot.JW}: \texttt{condExp}, \texttt{condExp\_\allowbreak{}isLinearMap}, \texttt{condExp\_\allowbreak{}one}, \texttt{condExp\_\allowbreak{}posSemidef} and 7 more; \texttt{Locality.LiebRobinson}: \texttt{measurable\_\allowbreak{}condExp}, \texttt{measurable\_\allowbreak{}condExp\_\allowbreak{}of\_\allowbreak{}continuous} \\
Theorem~5.2 & Lieb--Robinson bound for even interactions & \texttt{Locality.LiebRobinson}: \texttt{lieb\_\allowbreak{}robinson}, \texttt{lieb\_\allowbreak{}robinson\_\allowbreak{}of\_\allowbreak{}disjoint}, \texttt{lieb\_\allowbreak{}robinson\_\allowbreak{}of\_\allowbreak{}terms\_\allowbreak{}zero}, \texttt{chainWeight} and 3 more \\
Theorem~5.3 & propagation from a set & \texttt{Locality.Propagation}: \texttt{ck}, \texttt{rj}, \texttt{C1}, \texttt{C1\_\allowbreak{}le} and 7 more \\
Lemma~5.5 & propagation of almost-supported operators & \texttt{Locality.Propagation}: \texttt{setNorm\_\allowbreak{}dyn\_\allowbreak{}le} \\
Lemma~5.6 & time filters & \texttt{Locality.Filters}: \texttt{moment}, \texttt{HasFiniteMoments}, \texttt{setNorm\_\allowbreak{}timeFilter\_\allowbreak{}le}, \texttt{almostSupported\_\allowbreak{}timeFilter} \\
Lemma~5.7 & commutator with the Hamiltonian & \texttt{Locality.Filters}: \texttt{norm\_\allowbreak{}comm\_\allowbreak{}HN\_\allowbreak{}le}, \texttt{almostSupported\_\allowbreak{}comm\_\allowbreak{}HN} \\
Definition~5.8 & the filters & \texttt{Locality.Filters}: \texttt{IsCutoff}, \texttt{What}, \texttt{Wfilt}, \texttt{ClusterParams} and 3 more \\
Lemma~5.9 & the generator filter & \texttt{Locality.Filters}: \texttt{Wfilt\_\allowbreak{}im}, \texttt{Wfilt\_\allowbreak{}bounded}, \texttt{contDiffOn\_\allowbreak{}Wfilt}, \texttt{tendsto\_\allowbreak{}Wfilt\_\allowbreak{}zero} and 7 more \\
Lemma~5.10 & time representation of the clustering filter & \texttt{Locality.Filters}: \texttt{zetaHat}, \texttt{Fclust}, \texttt{cF}, \texttt{Gclust} and 6 more \\
Proposition~5.11 & locality of the clustering filter & \texttt{Locality.Filters}: \texttt{almostSupported\_\allowbreak{}clusterFilter} \\
Lemma~5.12 & locality calculus & \texttt{Locality.Calculus}: \texttt{norm\_\allowbreak{}comm\_\allowbreak{}le\_\allowbreak{}of\_\allowbreak{}disjoint}, \texttt{almostSupported\_\allowbreak{}add\_\allowbreak{}star\_\allowbreak{}mul}, \texttt{ordExp\_\allowbreak{}almostSupported}, \texttt{exp\_\allowbreak{}Qband\_\allowbreak{}add\_\allowbreak{}almostSupported} \\
Lemma~5.13 & bookkeeping & \texttt{Locality.Calculus}: \texttt{bookkeeping\_\allowbreak{}comm}, \texttt{IsONegInf.sum\_\allowbreak{}mul\_\allowbreak{}mul} \\
Setting~6.1 &  & \texttt{Assembly.Setting}: \texttt{IxSetting}; index core: \texttt{IndexCore.Translation}, \texttt{IndexCore.Statement} \\
Lemma~6.2 & separations & \texttt{Assembly.Setting}: \texttt{sep\_\allowbreak{}cuts}, \texttt{sep\_\allowbreak{}quarters}, \texttt{sep\_\allowbreak{}cut\_\allowbreak{}mid}, \texttt{sep\_\allowbreak{}cut\_\allowbreak{}opposite} and 2 more \\
Assumption~6.3 & gap in a sector & \texttt{Assembly.Setting}: \texttt{GapAssumption}, \texttt{ArgData.rank\_\allowbreak{}P}, \texttt{ArgData.P\_\allowbreak{}props}; index core: \texttt{IndexCore.Statement}, \texttt{IndexCore.Translation} \\
Assumption~6.4 & charge transport & \texttt{Assembly.Translation}: \texttt{TransportAssumption}; index core: \texttt{IndexCore.Statement} \\
Definition~6.5 & current, generator, dressed charge & \texttt{Assembly.Setting}: \texttt{Jminus}, \texttt{Jplus}, \texttt{Kfull}, \texttt{Kminus} and 2 more; index core: \texttt{IndexCore.Statement}, \texttt{IndexCore.Translation} \\
Lemma~6.6 & the split of the current & \texttt{Assembly.Setting}: \texttt{Jminus\_\allowbreak{}add\_\allowbreak{}Jplus}, \texttt{Kfull\_\allowbreak{}eq\_\allowbreak{}add}, \texttt{J\_\allowbreak{}isSelfAdjoint}, \texttt{K\_\allowbreak{}isSelfAdjoint} and 4 more; index core: \texttt{IndexCore.KP}, \texttt{IndexCore.Statement} \\
Lemma~6.7 & the in-sector identity & \texttt{Assembly.Setting}: \texttt{K\_\allowbreak{}comm\_\allowbreak{}P}; index core: \texttt{IndexCore.KP}, \texttt{IndexCore.Statement} \\
Lemma~6.8 & strip localization of the dressing & \texttt{Assembly.Setting}: \texttt{norm\_\allowbreak{}K\_\allowbreak{}le}, \texttt{K\_\allowbreak{}almostSupported}, \texttt{norm\_\allowbreak{}KShift\_\allowbreak{}le} \\
Proposition~6.10 & in-sector clustering & \texttt{Assembly.Inputs}: \texttt{Qf\_\allowbreak{}mul\_\allowbreak{}P}, \texttt{P\_\allowbreak{}mul\_\allowbreak{}comm\_\allowbreak{}mul\_\allowbreak{}P}, \texttt{clustering\_\allowbreak{}isONegInf}; index core: \texttt{IndexCore.Clustering}, \texttt{IndexCore.Statement} \\
Definition~6.11 & the objects of the argument & \texttt{Assembly.Inputs}: \texttt{KUminus}, \texttt{KUplus}, \texttt{QbarMinus}, \texttt{QbarPlus} and 7 more \emph{(for $U=\Theta$)}; index core: \texttt{IndexCore.Statement}, \texttt{IndexCore.Commutation} \\
Lemma~6.12 & localization and size & \texttt{Assembly.Inputs}: \texttt{table\_\allowbreak{}Q}, \texttt{table\_\allowbreak{}T}, \texttt{table\_\allowbreak{}J}, \texttt{table\_\allowbreak{}K} and 4 more \emph{(for $U=\Theta$)}; index core: \texttt{IndexCore.Statement} \\
Lemma~6.13 & commutation relations & \texttt{Assembly.Inputs}: \texttt{sepPairs}, \texttt{exact\_\allowbreak{}relations}, \texttt{sep\_\allowbreak{}isONegInf}, \texttt{comm\_\allowbreak{}Qmid\_\allowbreak{}isONegInf} and 3 more \emph{(for $U=\Theta$)}; index core: \texttt{IndexCore.Translation}, \texttt{IndexCore.Commutation} \\
Lemma~6.14 & Lemma~4.2 of \cite {BBDF21} in the sector & index core: \texttt{IndexCore.BBDF42}, \texttt{IndexCore.Clustering} \\
Corollary~6.15 & the localized unitaries commute with $P$ & \texttt{Assembly.Inputs}: \texttt{clu\_\allowbreak{}isONegInf}; index core: \texttt{IndexCore.BBDF42}, \texttt{IndexCore.Statement} \\
Lemma~6.16 & the differential equation and its stability & index core: \texttt{IndexCore.ODE} \\
Lemma~6.17 & the determinant at $2\pi $ & index core: \texttt{IndexCore.Determinant}, \texttt{IndexCore.Det} \\
Theorem~6.18 & index theorem for lowest-Landau-level Hamiltonians on the torus & \texttt{Assembly.Inputs}: \texttt{index\_\allowbreak{}theorem\_\allowbreak{}\ensuremath{\Theta}} \emph{(for $U=\Theta$)}; \texttt{IndexCore.Main}: \texttt{index\_\allowbreak{}core}, \texttt{index\_\allowbreak{}core\_\allowbreak{}400} (finite-dimensional form, general U) \\
Lemma~7.1 & the magnetic translation transports one orbital & \texttt{Assembly.Translation}: \texttt{theta\_\allowbreak{}commute}, \texttt{theta\_\allowbreak{}transport}, \texttt{theta\_\allowbreak{}exp}, \texttt{theta\_\allowbreak{}trace} and 1 more; index core: \texttt{IndexCore.Translation}, \texttt{IndexCore.Ring} \\
Corollary~7.2 & the translation corollary & \texttt{Assembly.Translation}: \texttt{translation\_\allowbreak{}corollary}; \texttt{IndexCore.Translation}: \texttt{Orbital.translation\_\allowbreak{}corollary} (finite content) \\
Corollary~7.3 & half filling with zero spin polarization & \texttt{Assembly.Translation}: \texttt{four\_\allowbreak{}dvd\_\allowbreak{}of\_\allowbreak{}halfFilling}, \texttt{four\_\allowbreak{}dvd\_\allowbreak{}hollowCore}, \texttt{no\_\allowbreak{}gapped\_\allowbreak{}patch}; \texttt{IndexCore.Translation}: \texttt{Orbital.four\_\allowbreak{}dvd\_\allowbreak{}p} (finite content) \\
Corollary~7.4 & odd $N$ and spin one half & \texttt{Assembly.Translation}: \texttt{no\_\allowbreak{}gap\_\allowbreak{}odd}, \texttt{many\_\allowbreak{}low\_\allowbreak{}odd}; \texttt{IndexCore.Translation}: \texttt{Orbital.odd\_\allowbreak{}plus}, \texttt{Orbital.odd\_\allowbreak{}minus} (finite content) \\
Theorem~8.1 & no uniform gap above the zero modes & \texttt{Assembly.Main}: \texttt{rectTorus}, \texttt{evenEig}, \texttt{hollow\_\allowbreak{}core\_\allowbreak{}index\_\allowbreak{}four\_\allowbreak{}dvd}, \texttt{hollow\_\allowbreak{}core\_\allowbreak{}index} and 1 more \\
Corollary~8.2 & odd $N$ & \texttt{Assembly.Main}: \texttt{oddEigPlus}, \texttt{oddEigMinus}, \texttt{odd\_\allowbreak{}E3\_\allowbreak{}tendsto}, \texttt{odd\_\allowbreak{}many\_\allowbreak{}small} \\
Lemma~8.3 & exchange of the sides & \texttt{Assembly.Main}: \texttt{swapTorus}, \texttt{rotFun}, \texttt{FfR}, \texttt{isLLL\_\allowbreak{}rotFun} and 6 more \\
Corollary~8.4 & uniformity over aspect ratios & \texttt{Assembly.Main}: \texttt{tauSet}, \texttt{E11}, \texttt{uniform\_\allowbreak{}E11\_\allowbreak{}tendsto}, \texttt{uniform\_\allowbreak{}E11\_\allowbreak{}eventually} \\
Proposition~B.1 & disjoint supports {\cite [Proposition~2.1]{NSY18}} & \texttt{Locality.LiebRobinson}: \texttt{commute\_\allowbreak{}of\_\allowbreak{}mem\_\allowbreak{}CARmEven\_\allowbreak{}of\_\allowbreak{}mem\_\allowbreak{}CARm}, \texttt{anticomm\_\allowbreak{}of\_\allowbreak{}mem\_\allowbreak{}CARm\_\allowbreak{}of\_\allowbreak{}isOdd} \\
Lemma~B.2 & \cite [Lemma~3.2]{NSY18} & \texttt{Locality.LiebRobinson}: \texttt{norm\_\allowbreak{}le\_\allowbreak{}of\_\allowbreak{}hasDerivAt\_\allowbreak{}comm} \\
Lemma~B.3 & the one-step estimate {\cite [Lemma~3.3]{NSY18}} & \texttt{Locality.LiebRobinson}: \texttt{norm\_\allowbreak{}comm\_\allowbreak{}dyn\_\allowbreak{}le\_\allowbreak{}oneStep} \\
\end{longtable}
\end{footnotesize}


\paragraph{Where the Lean proof takes another route.} The Lean proofs follow this paper's proofs, except in the following
places, where they use a different argument; the statements are unaffected. Where this paper names a constant the Lean's is
inside it, with one exception: the constant $C^H_k$ of Proposition~\ref{prop:slotH}(ii), which the paper's proof gives as
$55\,C_k$ and the Lean obtains as $192\pi\,C_k$ (the statement is existential).
\begin{itemize}[leftmargin=*]
\item The commutation $[\Theta,H_N]=0$ is proved twice: as Lemma~\ref{lem:covariance}, by the paper's route, and, for the
statement of Section~\ref{sec:setting}, from the covariance of the orbitals through the coefficient identity
$V_{a+1,b+1,c+1,d+1}=V_{abcd}$.
\item The sector-preserving approximant of Lemma~\ref{lem:embedding} by a finite average over gauge phases, in place of the
integral.
\item Lemma~\ref{lem:jets} by a real-variable argument rather than a holomorphic one; Theorem~\ref{thm:localform}(a) by reducing
the matrix entries on the whole Fock space to a form on pair amplitudes, in place of Step~6's coefficient matrix on the
two-particle space; Lemma~\ref{lem:evaluation}(iv) by geometric majorants.
\item The Lieb--Robinson iteration of Appendix~\ref{app:LR} as an induction on a per-slot bound, without the chain sums;
Lemma~\ref{lem:almostinputs} by propagating one best approximant at order $k+2$ instead of the dyadic decomposition;
Step~4 of Theorem~\ref{thm:propagation}(a) by shells of width $\lambda$ with weights $e^{-n}$ in place of the planar integral,
which gives $\sum_{n\ge0}e^{-n}(n+3)^2\le26$, that is $6.5\pi r_j^2(1+C_1\abs{t'})^2$ where the paper has
$8\pi r_j^2(1+C_1\abs{t'})^2$; the master inequality is then stated with the paper's constant; Lemma~\ref{lem:calculus}(c) by an
energy method.
\item The completeness half of Lemma~\ref{lem:orbitals}(ii) with the $y$-Fourier coefficients defined as integrals and their
first-order equation obtained by differentiation under the integral sign and an integration by parts, where the paper
differentiates the series term by term.
\item Lemma~\ref{lem:KP} and Lemma~\ref{lem:theta}(d) by applying the index core's lemmas to instances whose unused fields
carry trivial data, so that neither statement acquires a hypothesis this paper's lacks.
\item Lemma~\ref{lem:split}(ii)--(iii): the far region bounded by its supremum times the volume of the half-torus, in place of
the one-dimensional integral.
\item Corollary~\ref{cor:classwide-odd}(ii) by a finite pigeonhole over the scales $\gamma/4^i$, in place of compactness along a
subsequence, with the dimension bound $\binom{2N}{(N-1)/2}\ge(N+1)/2$ where the paper uses $\ge2N$; in
Corollary~\ref{cor:classwide-odd}(i) the paper's two cases are folded into one, $p-4k\ne0$ or $p-4k=0$.
\item Lemma~\ref{lem:exchange}: the conjugation of $H_N$ by the exchange unitary through the local form of
Theorem~\ref{thm:localform}(a), $H_N=\int h_R\,\dd R$, in place of the change of basis in $V_{abcd}$; the equality of the
sector eigenvalues through the counting function $E\mapsto\#\{k:E_k\le E\}$.
\item Appendix~\ref{app:kernel}, the lower bound. The canonical components of Section~\ref{kz:sec:many} as a sum over tuples of
orbitals, with the Slater-determinant form as a theorem, and the step from a function of the canonical shape back to a Fock
vector, which this appendix leaves silent, as a theorem (\texttt{canon\_surj}, by interpolation in every variable); the zero set
of $\vartheta_1$ in Lemma~\ref{kz:lem:theta}(c) from Jacobi's product, without the argument principle; Mathlib's Weierstrass
$\wp$, which differs from this appendix's by an additive constant, so that the relative state $d^{(p)}_{(1,1)}$ of the Lean is
this appendix's plus a multiple of $f_5^{(p)}$ (the span, and every statement, is unchanged). Mathlib has no analysis of several
complex variables, so the explicit states are carried as holomorphic numerators $F(U)\,J(x)^2J(y)^2\det(\cdots)$, the entries
$\vartheta_1^2\wp_\varepsilon$ made entire by the one-variable removable singularity, in the class of continuous functions
entire in each variable separately; (C1) through the vanishing of the holomorphic $Q$ in $f=\vartheta_1(u_s-u_t)^2Q$, and
Lemma~\ref{kz:lem:A} (2)~$\Rightarrow$~(3) using (C1) only at order two on the diagonal; (C2) by this appendix's route at the
generic points of a divisor $x_a=y_k$ and extended to every point of it along a complex line inside the divisor, on which only
countably many points are non-generic, by the continuity of $\partial_vf$.
\end{itemize}

\paragraph{Conventions.} The eigenvalues $E_k(N)$ of $H_N$ in a sector are those of Mathlib's spectral theorem for the
compression, the index clamped to $1\le k\le\dim$; Definition~\ref{def:bbdfobjects} and Lemmas~\ref{lem:objects}
and~\ref{lem:commutation} on the Fock space, and Theorem~\ref{thm:index}, are formalized for $U=\Theta$ only, the case this
paper's numbered results use (the index core's Lemma~\ref{lem:commutation}(i) and (iii) are for a general $U$);
$\varepsilon(N)$ is the largest of the twenty input quantities (for $U=\Theta$, (U2) and (U4) hold exactly and the twenty-one
separated commutators reduce to sixteen pairs, which with the four clustering commutators make twenty) and $M(N)=(1+4\pi C_h\norm{W_\gamma}_1)N^2$; the set norm
of Definition~\ref{def:almostsupported} and the distances
to sets take values in $[0,\infty]$. Two sentences of this paper were made precise by the formalization and read accordingly
in this version: Lemma~\ref{lem:almostinputs} and Lemma~\ref{lem:filters} state the hypothesis $\norm A_{X;\max\{k,3\}}<\infty$
that their proofs assumed (without it, with $X=\emptyset$ and the convention $0\cdot\infty=0$, their literal transcription is
false), and the sentences on odd $N$ in Section~\ref{sec:setting} and Corollary~\ref{cor:oddN} say $N\ge3$, since at $N=1$ the
sectors in question are two-dimensional and have no $E_3$. Some Lean statements carry hypotheses that this paper's statements
carry and their proofs do not use; they are kept, and \texttt{lean/README.md} lists them.
